% section_5_conclusion.tex — Discussion + Conclusion (merged per skeleton v5.1)
% Drafted per: docs/Draft/THESIS_SKELETON.md v5.1 §V (post-capex-drop 2026-04-22):
%   Summary, channel-asymmetry synthesis (answers 3 deferrals from §3.2/§4.1/§4.2),
%   mechanism disclosure, limitations, future research.

\section{Discussion and Conclusion}
\label{sec:conclusion}

\subsection{Summary of Findings}

The Cash Holdings hypothesis (HC, \S\ref{sec:main:hc}) is supported on the spontaneous Q\&A channel of CEO speech across all twelve specifications of suite H1, with contemporaneous and one-quarter-lead positive coefficients on UncAnsCEO that survive the four-step fixed-effects ladder; the Presentation channel of CEO speech (UncPreCEO) is null in the same panel. The Financial Constraint moderation hypothesis (HFC, \S\ref{sec:main:hfc}) confirms HC for the Q\&A channel and concentrates the precautionary cash response in Unrated firms relative to the rated pool: the UncAnsCEO\_c~$\times$~Unrated interaction is positive and significant in six of eight interaction cells, with all four lead-DV cells reaching $p<0.10$. The driver regressions (\S\ref{sec:additional:drivers}) confirm that the speech-uncertainty metric responds positively to the firm-level PRisk index of \citeA{hassan2020} across both CEO channels, and to the news-based macro indices of \citeA{baker2016} and \citeA{davis2016} preferentially through the Presentation channel. The outside-world reaction regressions (\S\ref{sec:additional:reaction}) show that the post-call $26$-day closing-quote relative bid-ask spread (Spread\textsubscript{25D}) widens contemporaneously with UncPreCEO and at a one-quarter lag with UncAnsCEO. The economic magnitude of the cash response is modest---a one-standard-deviation increase in UncAnsCEO corresponds to between $0.4\%$ and $0.7\%$ of the sample-mean cash-to-assets ratio across the FE ladder---but statistically robust, and the moderator specification confirms that UncPreCEO\_c and its Unrated interaction are null or marginal across the same ladder.

\subsection{Channel-Asymmetry Synthesis}

Three empirical patterns identify a Q\&A-versus-Presentation asymmetry across the analyses. First, the cash response of \S\ref{sec:main:hc} and \S\ref{sec:main:hfc} loads on the spontaneous Q\&A channel and not on the prepared Presentation channel. Second, the macro-uncertainty drivers of \S\ref{sec:additional:drivers} (US EPU, Global EPU) load preferentially on the Presentation channel and load the Q\&A channel only weakly under firm fixed effects. Third, the post-call spread response of \S\ref{sec:additional:reaction} is contemporaneous on the Presentation channel and lagged on the Q\&A channel. The firm-call-level PRisk driver of \S\ref{sec:additional:drivers} is a boundary condition: it loads both CEO channels at approximately equal magnitude and at high precision, which is descriptively consistent with the channel asymmetry being a function of the aggregation level of the regressor (firm-quarter versus monthly macroeconomic) rather than of the segment alone. Per the pre-commitment of \S\ref{sec:framework:precommit}, we report this asymmetry descriptively and do not theorize a structural mechanism for the prepared-versus-spontaneous split beyond the observation that the two segments carry distinct uncertainty content propagating through different downstream margins.

\subsection{Mechanism Disclosure}

The findings of \S\ref{sec:main} and \S\ref{sec:additional:reaction} operate on theoretically distinct margins. The cash-holdings results test the firm's financing margin: greater perceived future-funding uncertainty raises the precautionary liquidity buffer, consistent with the \citeA{opler1999} prediction extended to a within-firm panel resolution at quarterly frequency. The post-call spread results test a separate, market-side margin: liquidity providers widen quoted spreads when CEO speech raises perceived informed-trading risk, a canonical information-asymmetry mechanism distinct from the firm's own funding response. The present paper does not test a causal bridge linking the two margins---e.g., whether the spread widening drives the cash accumulation, or whether the cash accumulation calms a subsequent spread, or whether both respond independently to a common shock. The two findings are descriptively complementary; their joint identification is left to future work.

\subsection{Limitations}

Five caveats bound the interpretation. First, all analyses are observational. The speech-uncertainty regressors are not exogenously assigned, and the \S\ref{sec:main} and \S\ref{sec:additional:reaction} estimates identify associations consistent with the precautionary and information-asymmetry mechanisms invoked above without identifying either causally. Second, CEO speaker identification depends on the Capital~IQ structured speaker-metadata fields documented in \S\ref{sec:framework:measurement}; the CEO-partitioned sample covers approximately fifteen percentage points fewer calls than the manager-pool aggregate over the same period, with this coverage gap dominated by calls in which no speaker is tagged with a CEO role and a smaller residual of unresolved speaker-type placeholders. Third, the HC sample covers 2002--2018 and the HFC moderator analysis is bounded at 2016 by Compustat credit-rating coverage; HFC results in the 2017--2018 sub-period are therefore not tested. Fourth, the post-call spread window of \S\ref{sec:additional:reaction} is a 26-day construction (the call day plus 25 following trading days) chosen to capture market reaction sustained beyond the announcement window; alternative window lengths are not estimated in the body. Fifth, the data pipeline retains the internal column name \texttt{BGTLevel\_Spread} for backward compatibility with prior runs; the user-facing label Spread\textsubscript{25D} is used consistently in the body and tables, but legacy specification JSON and panel artifacts reflect the prior nomenclature.

\subsection{Future Research}

Five extensions are flagged as natural follow-ups. First, an identification strategy for the speech-uncertainty effect on cash, exploiting plausible exogenous shocks to CEO communication style (e.g., disclosure-rule changes, exogenous CEO turnover, or natural experiments in mandatory-disclosure timing) to recover causal estimates against the observational baseline of \S\ref{sec:main}. Second, a within-firm temporal-sequencing test linking the cash response of \S\ref{sec:main:hc} and the spread response of \S\ref{sec:additional:reaction} within a common firm-quarter window, to test whether the two margins are temporally ordered or operate in parallel. Third, additional cash-response moderators beyond the Unrated category of \S\ref{sec:main:hfc}, including size quartiles, dividend-payer status, and continuous credit-constraint indices, to identify complementary constraint dimensions that the binary FP scheme aggregates. Fourth, an OPSW log-cash-DV robustness specification of HC that estimates the cash response under the alternative log-form dependent variable, complementing the linear \citeA{bates2009} specification adopted in the body. Fifth, full pipeline rename of the spread variable to Spread\textsubscript{25D} across the underlying panel, builders, and specification JSON, removing the BGT-window legacy column name retained in this draft for backward compatibility.

\subsection{Concluding Remark}

The CEO partition of the call-level speech-uncertainty measure identifies a firm-quarter precautionary signal that the manager-pooled aggregate masks: the Q\&A channel of CEO speech is positively associated with contemporaneous and one-quarter-lead cash holdings (HC), with the response sharpest at credit-constrained Unrated firms (HFC). Construct validation against firm-level (\citeA{hassan2020}) and macro-level (\citeA{baker2016}, \citeA{davis2016}) uncertainty drivers confirms that the metric carries economically meaningful uncertainty content across multiple aggregation levels. Outside-world reaction evidence shows that market liquidity providers price both CEO speech segments into post-call quoted spreads, with the Presentation channel registering contemporaneously and the Q\&A channel registering with a one-quarter lag. The two channel asymmetries documented in the body---Q\&A dominant for cash and lagged spread, Presentation dominant for macro-driven speech-content variation and contemporaneous spread---are descriptively consistent with the prepared and spontaneous segments of an earnings call carrying distinct, economically meaningful uncertainty content that propagates through different downstream margins. The systematic application of the CEO partition to firm-quarter financing and external-reaction outcomes is the principal measurement contribution of this paper, and motivates the extensions outlined above.
